%
% Configuratie
%

% Preambule met standaardinstellingen
\documentclass[a4paper,oneside,11pt,final]{memoir}

% Gesplitst document
\usepackage{subfiles}

% Noot: zorg ervoor dat Nederlandse woord-splitsing geactiveerd is.
\usepackage[dutch]{babel}

% UTF8 gebruiken voor gebruik van alle symbolen
\usepackage[utf8]{inputenc}
\usepackage{eurosym}
\DeclareUnicodeCharacter{03A9}{\Omega}

% Noot: je kan het graphicxpakket een optie dvips of pdftex doorgeven
% in dat geval moet je ze ook aan iiiscriptie doorgeven, dus bijvoorbeeld
% \usepackage[dvips]{graphicx}
% \usepackage[dvips]{iiiscriptie}
\usepackage{graphicx}
\usepackage{iiiscriptie}

% Tabellen eleganter maken
\usepackage{booktabs}

% Navigeerbaarheid van hyperlinks in PDF
\usepackage{hyperref}

% Afkortingen intelligent gebruiken
\usepackage[printonlyused,withpage]{acronym}

% Math support
\usepackage{amssymb,amsmath}

% Extra functies
% Verkleinde margin entry
\setlength{\marginparwidth}{1.2in}
\let\oldmarginpar\marginpar
\renewcommand\marginpar[1] {\-\oldmarginpar[\raggedleft\footnotesize #1]%
{\raggedright\footnotesize #1}}

% Een TODO-entry
\newcommand{\todo}[1] {
	\addcontentsline{tdo}{todo}{\protect{#1}}
	\marginpar{#1}
}

% Hyperlink maken en URL in footnote tonen
\usepackage{hyperref}
\newcommand{\makeurl}[2]{\href{#1}{#2} \footnote{#1}}

% Compacte enumeraties
\newenvironment{enumerate_compact}{
\begin{enumerate}
  \setlength{\itemsep}{1pt}
  \setlength{\parskip}{0pt}
  \setlength{\parsep}{0pt}
}{\end{enumerate}}
\newenvironment{itemize_compact}{
\begin{itemize}
  \setlength{\itemsep}{1pt}
  \setlength{\parskip}{0pt}
  \setlength{\parsep}{0pt}
}{\end{itemize}}

% Float voor codefragmenten
\usepackage{float}
\floatstyle{ruled}
\newfloat{code}{thp}{lop}
\floatname{code}{Codefragment}


%
% Titelpagina
%

% Invullen velden
\departement{Departement Toegepaste Ingenieurswetenschappen}
\deptadres{Valentin Vaerwyckweg 1, 9000 Gent}
\studiejaar{1e Master Informatica}
\soortrapport{
Scriptie voorgedragen tot het behalen van de diploma van\\
MASTER IN DE TOEGEPASTE INGENIEURSWETENSCHAPPEN: INFORMATICA\\
}
\title{Ontwikkeling van een multimediaframework}
\bedrijfslogo{\includegraphics[height=18mm]{logomira}}
\author{Tim BESARD}
\promotoren{
Intern:& Leen BROUNS \\
Extern:& Philippe MOLLET}


%
% Inhoud
%

\begin{document}

\pagenumbering{roman}
%
% Titelpagina
%

\maketitle
\pagenumbering{roman}


%
% Abstract
%

\chapter*{Abstract}
\addcontentsline{toc}{chapter}{Abstract}
\label{chap:abstract}

%
% Voorwoord
%

\chapter*{Voorwoord}
\addcontentsline{toc}{chapter}{Voorwoord}
\label{chap:voorwoord}


%
% Inhoudstafel
%

\setlength\cftpartnumwidth{2em}

\newpage

\label{chap:inhoudstafel}
\tableofcontents

\newpage

\pagenumbering{arabic}


%
% Lijst met afbeeldingen
%

\listoffigures


%
% Lijst van afkortingen
%

\chapter*{Lijst van afkortingen}
\addcontentsline{toc}{chapter}{Lijst met afkortingen}
\label{chap:afkortingen}

\begin{acronym}[WYSIWYG]	% langste afkorting

\acro{upnp}[UPnP]{Universal Plug and Play}
\acro{mdns}[mDNS]{Multicast DNS}
\acro{dns}[DNS]{Domain Name System}
\acro{ssdp}[SSDP]{Simple Service Discovery Protocol}
\acro{dcp}[DCP]{Device Control Protocols}
\acro{ssl}[SSL]{Secure Socket Layer}
\acro{rpc}[RPC]{Remote Procedure Call}
\acro{webdav}[WebDAV]{Web-based Distributed Authoring and Versioning}
\acro{http}[HTTP]{Hypertext Transport Protocol}
\acro{rmi}[RMI]{Remote Method Invocation}
\acro{corba}[CORBA]{Common Object Request Broker Architecture}
\acro{soap}[SOAP]{Simple Object Access Protocol}
\acro{wysiwyg}[WYSIWYG]{What You See Is What You Get}
\acro{html}[HTML]{Hypertext Markup Language}
\acro{rest}[REST]{Representational State Transfer}
\acro{jvm}[JVM]{Java Virtual Machine}
\acro{svn}[SVN]{Subversion}

\end{acronym}



\newpage

\pagenumbering{arabic}
\subfile{1_introductie.tex}
\subfile{2_ontwerp.tex}
\subfile{3_server.tex}
\subfile{4_kiosk.tex}
\subfile{5_interface.tex}
\subfile{6_invoering.tex}

% Bibliografie
\bibliographystyle{plainnat}
\bibliography{verslag}


\end{document}
